%; whizzy paragraph -pdf xpdf -latex ./whizzypdfptex.sh
%; whizzy-paragraph "^\\\\begin{frame}"
% latex beamer presentation.
% platex, latex-beamer でコンパイルすることを想定。 

%     Tokyo Debian Meeting resources
%     Copyright (C) 2009 Junichi Uekawa
%     Copyright (C) 2009 Nobuhiro Iwamatsu

%     This program is free software; you can redistribute it and/or modify
%     it under the terms of the GNU General Public License as published by
%     the Free Software Foundation; either version 2 of the License, or
%     (at your option) any later version.

%     This program is distributed in the hope that it will be useful,
%     but WITHOUT ANY WARRANTY; without even the implied warreanty of
%     MERCHANTABILITY or FITNESS FOR A PARTICULAR PURPOSE.  See the
%     GNU General Public License for more details.

%     You should have received a copy of the GNU General Public License
%     along with this program; if not, write to the Free Software
%     Foundation, Inc., 51 Franklin St, Fifth Floor, Boston, MA  02110-1301 USA

\documentclass[cjk,dvipdfmx,12pt]{beamer}
\usetheme{Tokyo}
\usepackage{monthlypresentation}

%  preview (shell-command (concat "evince " (replace-regexp-in-string
%  "tex$" "pdf"(buffer-file-name)) "&")) 
%  presentation (shell-command (concat "xpdf -fullscreen " (replace-regexp-in-string "tex$" "pdf"(buffer-file-name)) "&"))
%  presentation (shell-command (concat "evince " (replace-regexp-in-string "tex$" "pdf"(buffer-file-name)) "&"))

%http://www.naney.org/diki/dk/hyperref.html
%日本語EUC系環境の時
\AtBeginDvi{\special{pdf:tounicode EUC-UCS2}}
%シフトJIS系環境の時
%\AtBeginDvi{\special{pdf:tounicode 90ms-RKSJ-UCS2}}

\newenvironment{commandlinesmall}%
{\VerbatimEnvironment
  \begin{Sbox}\begin{minipage}{1.0\hsize}\begin{fontsize}{8}{8} \begin{BVerbatim}}%
{\end{BVerbatim}\end{fontsize}\end{minipage}\end{Sbox}
  \setlength{\fboxsep}{8pt}
% start on a new paragraph

\vspace{6pt}% skip before
\fcolorbox{dancerdarkblue}{dancerlightblue}{\TheSbox}

\vspace{6pt}% skip after
}
%end of commandlinesmall

\title{第173回東京エリアDebian勉強会 \\　\\grml-debootstrapを用いた \\USB起動メモリの作成}
\subtitle{}
\author{NOKUBI Takatsugu (野首 貴嗣) \\knok@debian.org}
\date{2019-04-20}
\logo{\includegraphics[width=8cm]{image200607/openlogo-light.eps}}

\begin{document}

\section{表紙}

\begin{frame}
  \titlepage{}
\end{frame}


\section{自己紹介}

\begin{frame}{自己紹介}
  \begin{itemize}
  \item NOKUBI Takatsugu (野首 貴嗣)
  \item knok@debian.org / knok@daionet.gr.jp
  \item Twitter: @knok
  \item Debian developer since 2002
  \item bo 時代からのマシンを運用し続けている
  \end{itemize}
\end{frame}


\section{目次}

\begin{frame}{アジェンダ}
  \begin{itemize}
  \item USBメモリを使う動機
  \item grml-debootstrap概要
  \item セットアップ準備
  \item 設定
  \item ブートローダーの設定
  \end{itemize}
\end{frame}


\section{USBメモリを使う動機}
\emtext{USBメモリを使う動機}

\begin{frame}[containsverbatim]{クリーンな環境を構築できる}
  \begin{itemize}
  \item 内臓ストレージに影響なし
  \item USBメモリ以外の手法
    \begin{itemize}
    \item chroot環境(debootstrap)
      \begin{itemize}
      \item daemon類の扱いが面倒
      \end{itemize}
    \item コンテナ(LXC/LXD/Docker等)
    \item 完全仮想環境(VirtualBox, KVM等)
      \begin{itemize}
      \item 物理デバイスを扱いたい時大変
      \end{itemize}
    \end{itemize}
  \end{itemize}
\end{frame}


\begin{frame}[containsverbatim]{ポータブルな環境}
  \begin{itemize}
  \item 物理マシンの環境を汚さない
    \begin{itemize}
    \item 大抵はライセンスされたproprietary OSが入っている
    \item リカバリー領域、診断ツールなど
    \end{itemize}
  \item 一つのメディアを複数のデバイスで使える
  \item 複数のデバイスを用意すれば異なる環境を用意できる
  \item 物理デバイスすべてを扱える
    \begin{itemize}
    \item 仮想環境と比較して
    \end{itemize}
  \item NASの起動デバイスとして(Microserver)
  \end{itemize}
\end{frame}

\section{grml-debootstrap}
\emtext{grml-debootstrap}

\begin{frame}[containsverbatim]{grml-debootstrap}

GRML Live Linuxの一部
\begin{itemize}
  \item \url{http://grml.org/}
    \begin{itemize}
    \item システム管理者向けLive Linux環境
    \item FAIに対応
    \end{itemize}
  \item パッケージメンテナンス: Grml Team
  \item 関連ツール: grml-rescueboot, grml2usb
  \item ソースコード: \url{https://github.com/grml/grml-debootstrap}
\end{itemize}

\end{frame}


\section{セットアップ準備}
\emtext{セットアップ準備}

\subsection{USBメモリの選定}
\begin{frame}[containsverbatim]{選定基準}
  \begin{itemize}
  \item 用途と容量
    \begin{itemize}
    \item デスクトップ用途なら32GBもあれば十分
    \item ストレージを使わないサーバー用途なら16GB程度で十分
    \end{itemize}
  \item 速度
  \item 耐久性
    \begin{itemize}
    \item 安物はヘビーなアクセスが蓄積するとすぐ壊れる
    \end{itemize}
  \end{itemize}
\end{frame}

\begin{frame}[containsverbatim]{パーティショニング}
  \begin{itemize}
  \item EFIに対応する
    \begin{itemize}
    \item EFI領域を分ける必要あり
    \item GPTを使う
    \item オプション: MBRにも対応するなら2048セクタ目からパーティションを開始する
    \end{itemize}
  \item 方針: スワップパーティションは作らない
  \item rootファイルシステムの選択
    \begin{itemize}
    \item ext4
    \item f2fs - 最新のgrubでしか対応していない点に注意
    \end{itemize}
  \end{itemize}
\end{frame}

\begin{frame}[containsverbatim]{実行例}
  \begin{commandlinesmall}
$ sudo gdisk /dev/sda
GPT fdisk (gdisk) version 1.0.1
Command (? for help): n
Partition number (1-128, default 1): 
First sector (34-31358942, default = 2048) or {+-}size{KMGTP}: 
Last sector (2048-31358942, default = 31358942) or {+-}size{KMGTP}: +256M
Current type is 'Linux filesystem'
Hex code or GUID (L to show codes, Enter = 8300): ef00
Changed type of partition to 'EFI System'

Command (? for help): n
Partition number (2-128, default 2): 
First sector (34-31358942, default = 526336) or {+-}size{KMGTP}: 
Last sector (526336-31358942, default = 31358942) or {+-}size{KMGTP}: 
Current type is 'Linux filesystem'
Hex code or GUID (L to show codes, Enter = 8300): 
Changed type of partition to 'Linux filesystem'

Command (? for help): w

Final checks complete. About to write GPT data. THIS WILL OVERWRITE EXISTING
PARTITIONS!!

Do you want to proceed? (Y/N): y
OK; writing new GUID partition table (GPT) to /dev/sda.
The operation has completed successfully.
\end{commandlinesmall}
\end{frame}


\subsection{インストール}
\emtext{インストール}

\begin{frame}[containsverbatim]{ファイルシステム作成}
  \begin{commandlinesmall}
$ sudo mkfs.vfat -F32 /dev/sda1         
mkfs.fat 4.1 (2017-01-24)
sudo mkfs -t f2fs /dev/sda2

        F2FS-tools: mkfs.f2fs Ver: 1.7.0 (2016-07-28)

Info: Debug level = 0
Info: Label = 
Info: Trim is enabled
Info: Segments per section = 1
Info: Sections per zone = 1
Info: sector size = 512
Info: total sectors = 30832607 (15054 MB)
Info: zone aligned segment0 blkaddr: 256
Info: format version with
  "Linux version 4.18.0-0.bpo.1-amd64 (debian-kernel@lists.debian.org) (gcc version 6.3.0 20170516 (Debian 6.3.0-18+deb9u1)) #1 SMP Debian 4.18.6-1~bpo9+1 (2018-09-13)"
Info: Discarding device
Info: This device doesn't support BLKSECDISCARD
Info: This device doesn't support BLKDISCARD
Info: Overprovision ratio = 1.640%
Info: Overprovision segments = 249 (GC reserved = 129)
Info: format successful
  \end{commandlinesmall}
\end{frame}

\begin{frame}[containsverbatim]{マウントとセットアップ}
--efiオプションはEFI環境で起動していないと動作しない点に注意
  \begin{commandlinesmall}
$ sudo grml-debootstrap -m http://deb.debian.org/debian -r stretch -t /dev/sda2 --efi /dev/sda1 --grub /dev/sda --arch amd64 --filesystem ext4 --password toor
 * EFI support detected.
 * grml-debootstrap [0.78] - Please recheck configuration before execution:

   Target:          /dev/sda2
   Install grub:    /dev/sda
   Install efi:     /dev/sda1
   Using release:   stretch
   Using hostname:  xps13
   Using mirror:    http://deb.debian.org/debian
   Using arch:      amd64
   Config files:    /etc/debootstrap

   Important! Continuing will delete all data from /dev/sda2!

 * Is this ok for you? [y/N] y
 * EFI partition /dev/sda1 seems to have a FAT filesystem, not modifying.
 * Running mkfs.f2fs  on /dev/sda2

        F2FS-tools: mkfs.f2fs Ver: 1.7.0 (2016-07-28)

Info: Debug level = 0
Info: Label =
Info: Trim is enabled
Info: Segments per section = 1
Info: Sections per zone = 1
Info: sector size = 512
Info: total sectors = 30832607 (15054 MB)
Info: zone aligned segment0 blkaddr: 256
Info: format version with
  "Linux version 4.18.0-0.bpo.1-amd64 (debian-kernel@lists.debian.org) (gcc version 6.3.0 20170516 (Debian 6.3.0-18+deb9u1)) #1 SMP Debian 4.18.6-1~bpo9+1 (2018-09-13)"
Info: Discarding device
Info: This device doesn't support BLKSECDISCARD
Info: This device doesn't support BLKDISCARD
Info: Overprovision ratio = 1.640%
Info: Overprovision segments = 249 (GC reserved = 129)
Info: format successful
blockdev: ioctl error on BLKRRPART: Device or resource busy
/usr/sbin/grml-debootstrap: line 1119: filesystem: command not found
 * Mounting /dev/sda2 to /mnt/debootstrap.4583
 * Running debootstrap  for release stretch (amd64) using http://deb.debian.org/debian
 * Executing: debootstrap --arch amd64   stretch /mnt/debootstrap.4583 http://deb.debian.org/debian
I: Retrieving InRelease
I: Retrieving Release
I: Retrieving Release.gpg
I: Checking Release signature
I: Valid Release signature (key id 067E3C456BAE240ACEE88F6FEF0F382A1A7B6500)
I: Retrieving Packages
I: Validating Packages
I: Resolving dependencies of required packages...
I: Resolving dependencies of base packages...


Finished chroot installation, exiting.
 * Removing chroot-script again
 * Unmount /mnt/debootstrap.4583
 * Removing /var/cache/grml-debootstrap/variables_sda2
 * Removing /var/cache/grml-debootstrap/stages_sda2
 * Finished execution of grml-debootstrap. Enjoy your Debian system.
  \end{commandlinesmall}
\end{frame}

\begin{frame}[containsverbatim]{手作業}

EFIパーティションの/EFI/BOOT以下にBOOTX64.EFIという名称でファイルを配置
\begin{commandline}
# mount /dev/sda1 /mnt
# cd /mnt/EFI
# mkdir BOOT
# cp /boot/efi/EFI/debian/grubx64.efi BOOT/BOOTX64.EFI
# umount /mnt
\end{commandline}
同様にgrub-efi-ia32を用いてBOOT/BOOTx32.EFIを用意すれば32bit UEFI対応もできる
\\
MBRに書き込めばEFI/MBR両対応も可能
\end{frame}

\section{動作確認}
\emtext{動作確認}

\begin{frame}[containsverbatim]{VirtualBoxによる動作確認}
  \begin{itemize}
  \item diskグループに追加
    \begin{itemize}
    \item sudo adduser user disk
    \item ログインし直す
    \end{itemize}
  \item Raw diskに対応するvmdkの作成
    \begin{commandlinesmall}
      $ vboxmanage internalcommands createrawvmdk -filename \
      ~/usb.vmdk -rawdisk /dev/sda
    \end{commandlinesmall}
  \item VirtualBoxで仮想マシンを作成
    \begin{itemize}
    \item ストレージに先程作成したusb.vmdkを指定
    \item UEFIを有効化
    \end{itemize}
  \item 起動
  \end{itemize}
\end{frame}

\begin{frame}[containsverbatim]{参考文献}
  \begin{itemize}
  \item 起動ストレージ
    \begin{itemize}
    \item fstab, grub.cfgにUUIDで記述
    \end{itemize}
  \item Network
    \begin{itemize}
    \item if addrでインターフェース名を確認してdhclient
    \end{itemize}
  \item Timezone
    \begin{itemize}
    \item dpkg-reconfigure tzdata
    \end{itemize}
  \item Hardware Clock
    \begin{itemize}
    \item /etc/adjtimeの末尾修正
    \item LOCALでRTCをローカル扱いに
    \end{itemize}
  \item Keymap
    \begin{itemize}
    \item dpkg-reconfigure keyboard-configuration
    \item \url{https://wiki.debian.org/Keyboard}
    \item お好みで XKBOPTIONS="ctrl:swapcaps"
    \end{itemize}
  \end{itemize}
\end{frame}

\begin{frame}[containsverbatim]{参考文献}
  \begin{itemize}
  \item grml-debootstrapを紹介
    \begin{itemize}
    \item インストーラを使わずストレージにDebian環境を構築
    \end{itemize}
  \item USBメモリを使う動機
    \begin{itemize}
    \item 複数のデバイスで同じ環境を利用できる
    \item メディアを分けることで異なる用途に使える
    \item 起動マシンの環境を汚さない
    \end{itemize}
  \item UEFI環境向けにデフォルトのブートローダーを配置
  \item VirtualBoxを動作テストに使える
    \begin{itemize}
    \item 細かい調整、パッケージの変更などもできる
    \end{itemize}
  \end{itemize}
\end{frame}

\section{参考文献}
\emtext{参考文献}

\begin{frame}[containsverbatim]{参考文献}

参考文献
  \begin{itemize}
  \item BIOS, UEFI両方で起動可能なdebianインストールUSBメディアを作る \url{http://pman0214.github.io/blog/debian-install-bios-efi.html}
  \item 32bit UEFI 搭載の2-in-1 PCにubuntu14.04をインストールした記録 \url{https://qiita.com/shirotamago/items/a4b0c8863a492abe50ad}
   \item Boot your USB Drive in VirtualBox \url{http://agnipulse.com/2009/07/boot-your-usb-drive-in-virtualbox/}
  \end{itemize}
\end{frame}

\begin{frame}[containsverbatim]{ファイルシステムについて}
  F2FSはどうか?
  \begin{itemize}
  \item 寿命は伸びるかも?(未検証)
    \begin{itemize}
    \item USBメモリでTRIMは発行できるか
    \item デバイス依存, hdparm -Iで確認
    \end{itemize}
  \item stretchのgrub2はf2fsをサポートしていない
    \begin{itemize}
      \item /boot を別に作ることで対応可能
    \end{itemize}
  \item パフォーマンス
    \begin{itemize}
    \item F2FS Benchmarks From USB Flash Storage (Phoronix) \url{https://www.phoronix.com/scan.php?page=article&item=linux_f2fs_usb3&num=1}
    \item 一長一短?
    \end{itemize}
  \end{itemize}
\end{frame}


\end{document}


;;; Local Variables: ***
;;; outline-regexp: "\\([ 	]*\\\\\\(documentstyle\\|documentclass\\|emtext\\|section\\|begin{frame}\\)\\*?[ 	]*[[{]\\|[]+\\)" ***
;;; End: ***
